Set \(d_\tau=c(\tau)-u\) and \(J=\sum_\tau d_\tau\), and write
\(\widetilde\phi_\vartheta\) for the lift before scaling. It is the
sum of the vacuum kernels over the classes of \(U(1)\), multiplied
by the level index on \(U(1)\) at the chosen places.
In the rank-one inner product formula the vacuum matrix-coefficient
sum away from the level multiplies the local masses at the level,
the archimedean integral and the quotient volumes. Let \(S\) consist
of the chosen and twist places and their conjugates. The vacuum sum is
\[
 \frac{L^S(u,(\nu\vartheta)^{-1})}{\zeta_{E^+}(2u)}
    \prod_{w\ \text{at\ level}}(1-q_w^{-2u})^{-1}.
\]
The product runs over the original places before adjoining their
conjugates. The sign factors at a place and its conjugate combine
to give the corresponding factor of \(\zeta_{E^+}(2u)\).

At a chosen place the column has one prescribed reduction, so its
mean square is \(q_w^{-2u}\) and the target level index is \(q_w^{2u}-1\).
Including the compensation for the omitted Euler factor gives \(1\).
At a twist place the selected coordinate is a unit in a fraction
\((q_w-1)/q_w\) of the columns, and its character has absolute value \(1\).
The index of the line stabilizer is \((q_w^{2u}-1)/(q_w-1)\), giving
\[
 \frac{q_w-1}{q_w}\,
 \frac{q_w^{2u}-1}{q_w-1}\,
 (1-q_w^{-2u})^{-1}=q_w^{2u-1}.
\]
The source level index is \(\prod_{w\ \mathrm{chosen}}(q_w-1)\).
It occurs once in the denominator of the block mass and twice in the
squared norm of the scaled sum, leaving one factor. At infinity the
Fock inner product contributes \(\prod_\tau d_\tau!\), and the
matrix-coefficient integral contributes the dimension of the Weyl
space \(\mathcal W\). Thus, writing \(A_u\) for the product of the
quotient volumes divided by \(\zeta_{E^+}(2u)\), we have
\[
 \|\widetilde\phi_\vartheta\|^2
 = A_u \prod_{w\ \mathrm{chosen}}(q_w-1)\,
     \prod_\tau d_\tau!\,\dim\mathcal W\,
     L^S(u,(\nu\vartheta)^{-1})\,N_\vartheta^{2u-1}.
\]
The vacuum and derivative scales passing to \(\phi_\vartheta\) are
independent of \(\vartheta\). They cancel in the ratio, and taking
absolute values gives \eqref{eq:main-pair-norm}.
\leanlinks{\leansrc{New/OneLift/HeckeValues/CharacterLiftNorm}{Local factors in the rank-one norm}
  \leanlinksep \leansrc{New/OneLift/HeckeValues/TwistedRatio}{Ratio of norms under twisting}}

Fix \(\sigma\in\operatorname{Aut}(\mathbb C/K)\) and conjugate the
untwisted squared norm. The two slots in the pairing involve
\(\sigma\) and \(\mathrm{conj}\circ\sigma\circ\mathrm{conj}\),
respectively. These induce the same permutation of the infinite
places of \(E\), so the orders \(d_\tau\) are permuted.
When the two conjugate lifts are written using the same finite
classes, the multipliers of the level character are inverse in
the pairing and cancel. The norm retains a common vacuum factor
and the derivative scale \(M^J\Omega^{-J}\). Consequently
\[
 \log|\sigma(\|\phi_1\|^2)|
   =c_{u,\sigma}+\sum_\tau s_\infty(d_\tau)+
      \log\dim\mathcal W^\sigma+\log|L^S(u,\nu_\sigma^{-1})|,
\]
where \(c_{u,\sigma}\) depends only on the rank and \(\sigma\).
Since \(u\ge2\), the Euler product bounds the absolute value of the
last logarithm by \(4f\). The dimension is at least \(1\), and it is
\(1\) when all derivative orders are zero. The last two terms
thus sum to at least \(-4f\), and to at most \(4f\) in degree zero.
\leanlinks{\leansrc{New/OneLift/HeckeValues/AlgebraicPartSize}{The two conjugations and derivative normalization}}
